\chapter{État de l'art de la menace logicielle et mécanismes d'évasion}

%===============================================================
\section*{Périmètre}
\addcontentsline{toc}{section}{Périmètre}

La souche analysée dans la suite de ce mémoire est un binaire ELF destiné à un environnement GNU/Linux. Cet état de l'art est donc construit autour des systèmes Linux, en particulier dans leurs usages serveur et cloud, et ne mobilise l'écosystème Windows qu'à titre de contrepoint lorsque la comparaison éclaire un mécanisme.

Ce choix n'est pas seulement dicté par l'objet d'étude. La littérature disponible est très déséquilibrée : les travaux sur la détection de binaires PE forment un corpus mature, doté de jeux de données publics et de protocoles d'évaluation partagés, tandis que la recherche équivalente sur les binaires ELF reste dispersée et beaucoup moins outillée \cite{crossArchSurvey}. Ce déséquilibre a une conséquence directe sur le travail présenté ici, puisqu'il interdit de reprendre telles quelles les méthodes d'évaluation du monde Windows, et il constitue en lui-même une partie de la justification du sujet.

Les sources mobilisées relèvent de trois registres qui n'ont pas la même valeur probante : travaux académiques évalués par les pairs, rapports d'investigation publiés par des équipes de réponse à incident, et documentation de référence des cadres d'analyse. Le texte signale explicitement les cas où une capacité relève de l'hypothèse plutôt que de l'observation.

%===============================================================
\section{Définition et taxonomie du malware moderne}

Un malware ne se définit plus par sa capacité de réplication, comme pouvaient l'être les premiers virus des années 1980. Ce qui le caractérise est l'intention : tout code conçu pour s'introduire dans un système, en perturber le fonctionnement, en extraire des données ou y maintenir un accès non autorisé, sans le consentement de son propriétaire.

Le critère est bien l'intention, et non la technique employée. Un outil de prise en main à distance est un logiciel d'administration parfaitement légitime quand un exploitant l'installe sur ses serveurs, et une porte dérobée quand un attaquant l'y dépose. Les deux peuvent être le même programme. Cette ambiguïté traverse tout ce chapitre : elle explique pourquoi la détection est difficile, et pourquoi elle ne se ramène jamais à reconnaître un objet dans une liste.

Le déplacement de définition accompagne un changement de nature de l'activité. La production de code malveillant, longtemps portée par des motivations de démonstration technique, s'est structurée en filière, avec une division du travail et des objectifs financiers ou stratégiques. Le volume de nouveaux échantillons recensés quotidiennement par les observatoires du secteur donne la mesure de cette industrialisation \cite{avtest}.

%---------------------------------------------------------------
\subsection{Nomenclature et standards d'identification}

Deux analystes qui examinent le même fichier dans deux entreprises différentes doivent pouvoir établir qu'ils parlent bien du même objet. C'est le rôle des conventions de nommage, dont la plus ancienne et la plus diffusée est celle du CARO, définie au début des années 1990 et que Microsoft déclare suivre pour ses propres détections \cite{msnaming}. Un identifiant y est composé de cinq champs, à la manière d'un nom d'espèce en biologie, du plus général au plus précis :

\begin{center}
\texttt{Type : Plateforme / Famille . Variante ! Suffixe}
\end{center}

Le tableau \ref{tab:nommage} décompose un identifiant réel selon ce schéma.

\begin{table}[h!]
\centering
\caption{Lecture de l'identifiant \texttt{Trojan:Linux/XorDDoS.A!ml}}
\label{tab:nommage}
\begin{tabular}{|l|l|p{7cm}|}
\hline
Champ & Valeur & Signification \\
\hline
Type & \texttt{Trojan} & Ce que le code fait : il se présente comme légitime pour être exécuté. Autres valeurs courantes : \texttt{Backdoor} (accès distant), \texttt{Ransom} (rançongiciel), \texttt{Worm} (propagation autonome), \texttt{PWS} (vol de mots de passe). \\
\hline
Plateforme & \texttt{Linux} & L'environnement visé. Peut aussi désigner un langage interprété (\texttt{Python}, \texttt{Shell}) ou une version de Windows. \\
\hline
Famille & \texttt{XorDDoS} & Le nom donné à la découverte, partagé par tous les échantillons dont le code est jugé assez proche pour supposer une origine commune. \\
\hline
Variante & \texttt{A} & Une lettre attribuée dans l'ordre des versions successives : \texttt{.AF} succède à \texttt{.AE}. \\
\hline
Suffixe & \texttt{!ml} & Une précision technique facultative. Ici, la détection provient d'un modèle d'apprentissage automatique et non d'une règle écrite à la main. \\
\hline
\end{tabular}
\end{table}

Sur les cibles Linux, le champ plateforme masque une difficulté propre au reste de ce mémoire. Un même code source y est fréquemment compilé pour plusieurs architectures matérielles, x86-64 pour les serveurs, ARM ou MIPS pour les équipements embarqués. Les binaires produits sont différents alors que le comportement est identique, et le champ plateforme n'en rend pas compte.

\paragraph{Les limites du schéma.} Présenter cette convention comme un standard sectoriel serait inexact. L'un de ses auteurs reconnaît qu'aucun produit n'en respecte intégralement les règles et que le CARO n'a jamais disposé de l'autorité nécessaire pour les imposer \cite{bontchev2005}. Deux éditeurs attribuent donc régulièrement des noms différents au même échantillon, et les tentatives ultérieures d'unification ont toutes échoué.

Le problème n'est pas seulement documentaire. Entraîner un modèle à reconnaître des familles suppose de disposer d'échantillons correctement étiquetés ; or l'étiquette de référence, en pratique, est le verdict des antivirus. Si ces verdicts se contredisent, les étiquettes sont bruitées, et le modèle apprend en partie ce désaccord. Des outils de consolidation existent, mais ils héritent des divergences qu'ils agrègent.

%---------------------------------------------------------------
\subsection{Tactique, technique, vecteur, capacité et charge utile}

Cinq notions se confondent couramment, et leur confusion fausse l'analyse de risque : on croit se protéger d'une menace quand on ne se protège que d'un de ses outils. Le cadre MITRE ATT\&CK \cite{attack} fournit une grille qui les sépare. Une comparaison avec un cambriolage en rend la logique immédiate, et le tableau \ref{tab:attack} la transpose sur un cas informatique.

\begin{table}[h!]
\centering
\caption{Les cinq niveaux d'analyse d'une intrusion}
\label{tab:attack}
\begin{tabular}{|l|p{3.6cm}|p{3.4cm}|p{4.2cm}|}
\hline
Niveau & Question posée & Cambriolage & Intrusion sur un serveur Linux \\
\hline
Tactique & Que cherche l'attaquant à cette étape ? & Entrer dans la maison & Obtenir un premier accès (TA0001) \\
\hline
Technique & Comment s'y prend-il ? & Crocheter la serrure & Exploiter une application exposée (T1190) \\
\hline
Vecteur & Par quel support concret ? & La porte de service & Un service web accessible depuis internet \\
\hline
Capacité & Que sait-il faire une fois entré ? & Ouvrir un coffre & Élever ses privilèges, se dissimuler \\
\hline
Charge utile & Qu'emporte-t-il au final ? & Les bijoux & Chiffrer les données, exfiltrer des clés \\
\hline
\end{tabular}
\end{table}

Une sixième notion traverse les autres : l'évasion défensive (TA0005), qui regroupe tout ce que l'attaquant met en \oe uvre pour ne pas être vu pendant qu'il agit. Elle constitue l'objet principal de ce chapitre, et son cambrioleur porterait des gants.

L'intérêt de cette grille est pratique. Une signature protège d'un échantillon, c'est-à-dire d'un cambrioleur précis dont on possède la photographie ; une détection comportementale protège d'une technique, donc de quiconque crochète une serrure. La première approche est exacte mais s'échappe dès que l'apparence change, la seconde tolère la variation mais produit des faux positifs, puisque le propriétaire aussi manipule sa serrure. Tout l'état de l'art de la détection tient dans cet arbitrage, et la section \ref{sec:detection} y revient.

%===============================================================
\section{Classification par objectifs opérationnels}

Classer les menaces par l'objectif de leurs auteurs est plus utile en pratique qu'une taxonomie par type de fichier.

Le rançongiciel chiffre les données et monnaye leur restitution. Sur Linux, il a suivi la valeur : les variantes ELF visant les hyperviseurs de virtualisation permettent de rendre indisponibles des dizaines de machines virtuelles par une seule exécution, ce qui déplace le rapport coût-bénéfice en faveur de l'attaquant.

Le voleur d'informations extrait discrètement des données valorisables. Sur un poste de travail il vise des cookies et des identifiants ; sur un serveur ou une machine de développement, il vise des secrets d'infrastructure, jetons d'API, clés SSH, variables d'environnement d'intégration continue, dont la compromission ouvre un accès bien plus large que la machine elle-même.

La porte dérobée assure un accès persistant à distance. C'est la catégorie la plus représentée dans les compromissions de serveurs, et celle qui concentre l'essentiel des techniques de dissimulation étudiées en section \ref{sec:runtime}.

Le mineur de cryptomonnaie détourne les ressources de calcul. Longtemps considéré comme une nuisance mineure, il sert en pratique de révélateur : sa présence signale une chaîne d'accès qui aurait tout aussi bien pu porter une charge destructive.

Le bot enrôle la machine dans un réseau contrôlé à distance. Les familles visant les équipements embarqués sous Linux, dérivées pour beaucoup d'un même code source diffusé publiquement, illustrent le problème de la compilation croisée : un même comportement se décline en dizaines de binaires dont la structure varie selon l'architecture cible.

Ces catégories ne sont pas étanches, et un échantillon combine ordinairement plusieurs fonctions.

%---------------------------------------------------------------
\subsection{L'écosystème économique}
\label{sec:economie}

La spécialisation des acteurs explique une part de cette polyvalence. Le modèle du
rançongiciel comme service sépare ceux qui écrivent le code de ceux qui conduisent
les intrusions. En amont, des courtiers revendent des accès déjà obtenus sans
participer à l'exploitation ; l'ANSSI en fait des intermédiaires devenus
indispensables à l'écosystème, en particulier pour les rançongiciels
\cite{anssiPanorama}.

Le mode opératoire Houken l'illustre. L'exploitation des vulnérabilités initiales y
est nettement plus soignée que la suite de la compromission, écart qui trahit le
passage de relais entre un courtier et les acteurs venus après lui. Autrement dit,
la charge finale ne renseigne plus sur le niveau technique de celui qui a ouvert la
porte.

Les volumes du CERT-FR donnent l'échelle. L'agence a connu 1~366 incidents sur
l'exercice 2025, contre 831 en 2022. Le rançongiciel recule légèrement, à 128
compromissions, et aucune souche ne domine : Qilin pèse 21~\% des cas, Akira 9~\%,
LockBit 5~\%, et plus d'une dizaine de souches inédites apparaissent sur au moins un
incident dans l'année. L'exfiltration de données, elle, passe de 130 à 196
incidents\cite{anssiPanorama}.

Le déplacement le plus gênant pour la détection est ailleurs. Plusieurs courtiers en
accès, dont TA577, TA571 et TA544, ont abandonné les chargeurs et les botnets au
profit d'outils d'administration à distance légitimes et de binaires natifs du
système. Lumma Stealer, de son côté, renouvelle ses domaines de commande et contrôle
plusieurs fois par semaine, un même échantillon pouvant en embarquer une dizaine.

Deux pressions pèsent donc sur une défense construite sur des artefacts figés. La
fragmentation d'abord : sans souche dominante, il faut couvrir une longue traîne de
familles pour un rendement décroissant. Le reste est plus radical. Quand
l'attaquant se sert d'un outil d'administration signé plutôt que d'un programme
malveillant, il n'y a plus d'artefact à reconnaître.

%===============================================================
\section{Évasion statique : rendre le binaire illisible}

L'analyse statique consiste à examiner un fichier sans l'exécuter. Les techniques qui la contrarient forment la première ligne d'évasion.

%---------------------------------------------------------------
\subsection{Réduction de la surface d'analyse}

Sur ELF, les mesures les plus efficaces sont souvent les plus élémentaires. La suppression des symboles prive l'analyste des noms de fonctions et de variables. La liaison statique fusionne les bibliothèques dans le binaire et supprime la table d'importation, qui constitue par ailleurs l'un des indicateurs les plus utilisés pour caractériser un exécutable. Le recours à des langages compilés modernes, Go ou Rust, produit des binaires volumineux dont l'exécution est dominée par le code de la bibliothèque standard, ce qui noie la logique métier dans un bruit considérable sans qu'aucune intention d'obfuscation ait été nécessaire.

L'altération volontaire des en-têtes et des tables de sections complète l'ensemble : le chargeur du noyau tolère des structures que les outils d'analyse refusent, et cet écart de tolérance suffit à mettre en échec un outillage automatisé.

%---------------------------------------------------------------
\subsection{Empaquetage}

Un packer compresse ou chiffre l'exécutable original et lui adjoint une routine de déballage exécutée au lancement. Sur ELF, l'usage d'un packer public dont les marqueurs ont été modifiés reste très répandu, précisément parce qu'il neutralise les outils de déballage automatique sans exiger de développement propre.

Une taxonomie des structures de déballage a été établie sur le monde PE, distinguant les modèles à couche unique des schémas incrémentaux où le code original n'est jamais entièrement reconstitué en mémoire à un instant donné \cite{ugartepedrero2015}. Cette dernière catégorie invalide l'hypothèse implicite de nombreuses chaînes d'analyse automatisée, celle d'un instant de capture où le code complet serait disponible.

%---------------------------------------------------------------
\subsection{Polymorphisme, métamorphisme et virtualisation}

Le polymorphisme chiffre le corps du code et fait varier la clé à chaque exécution : l'empreinte sur disque diffère, le contenu déchiffré en mémoire reste identique. Le métamorphisme va plus loin en réécrivant le code lui-même à chaque infection, par substitution d'instructions équivalentes, permutation de registres ou dispersion en fragments reliés par des sauts \cite{szor2005}.

Ces moteurs ont été développés et documentés dans le contexte des virus DOS et Windows. Ils sont nettement plus rares sur ELF, pour une raison économique : maintenir un moteur de réécriture correct sur plusieurs jeux d'instructions coûte cher, alors que l'écosystème Linux offre des voies d'évasion moins onéreuses, décrites en section \ref{sec:runtime}. L'observation mérite d'être retenue, car elle indique où porter l'effort défensif.

La virtualisation de code, qui traduit les instructions en un bytecode propriétaire interprété par une machine virtuelle embarquée, et les prédicats opaques, conditions dont le résultat est connu de l'auteur mais indécidable pour un désassembleur \cite{collberg1997}, complètent l'arsenal. Ces deux techniques restent surtout associées aux protecteurs commerciaux.

\begin{table}[h!]
\centering
\caption{Techniques d'évasion statique observées sur ELF}
\label{tab:evasion-statique}
\begin{tabular}{|l|p{5.5cm}|p{3.5cm}|}
\hline
Technique & Principe & Fréquence sur ELF \\
\hline
Suppression des symboles & Retrait des noms de fonctions et variables & Très courante \\
\hline
Liaison statique & Suppression de la table d'importation & Très courante \\
\hline
En-têtes altérés & Structures tolérées par le noyau, refusées par l'outillage & Courante \\
\hline
Empaquetage & Compression ou chiffrement, déballage au lancement & Courante \\
\hline
Déballage incrémental & Reconstitution partielle et successive en mémoire & Peu fréquente \\
\hline
Polymorphisme & Chiffrement du corps, clé variable & Peu fréquente \\
\hline
Métamorphisme & Réécriture du code à chaque infection & Rare (coût multi-architectures) \\
\hline
Virtualisation & Bytecode propriétaire et interpréteur embarqué & Rare \\
\hline
\end{tabular}
\end{table}

%===============================================================
\section{Évasion à l'exécution : échapper à l'observation}
\label{sec:runtime}

L'obfuscation traite le problème de la lecture du code. Elle ne dit rien du second front, celui de l'exécution sous surveillance, où se joue l'essentiel des compromissions de serveurs. Un code parfaitement obfusqué qui déclenche un chiffrement massif sera arrêté ; un code trivial qui n'émet aucun signal observable ne le sera pas.

%---------------------------------------------------------------
\subsection{Anti-analyse et conditionnement de l'environnement}

Avant de déployer sa charge, un code cherche à établir s'il s'exécute sur une cible réelle ou dans un environnement d'analyse. Sous Linux, les vérifications les plus simples passent par le pseudo-système de fichiers \texttt{/proc}, qui expose l'état du processus, y compris la présence d'un débogueur attaché. L'auto-attachement par \texttt{ptrace} produit le même résultat par un autre chemin : un processus ne pouvant être tracé qu'une fois, réussir à se tracer soi-même prouve l'absence d'observateur.

D'autres contrôles portent sur la nature de l'hôte, avec la recherche d'indices de virtualisation ou de conteneurisation, sur la temporisation, une exécution instrumentée étant plus lente d'un facteur mesurable, ou sur des caractéristiques d'environnement propres à la cible visée. Le conditionnement environnemental (T1480.001) pousse la logique jusqu'au bout : la clé de déchiffrement est dérivée de propriétés de la machine ciblée, de sorte que le code ne se déchiffre nulle part ailleurs, y compris entre les mains du défenseur.

%---------------------------------------------------------------
\subsection{Exécution sans fichier}

L'analyse statique suppose un fichier à analyser. L'appel système \texttt{memfd\_create} crée un objet mémoire anonyme utilisable comme descripteur exécutable, ce qui permet de charger et d'exécuter un binaire sans jamais l'écrire sur le système de fichiers. Le montage en mémoire vive de \texttt{/dev/shm} offre une variante moins discrète mais tout aussi accessible.

Le corollaire est important pour la démarche défensive : sur un hôte ainsi compromis, il n'existe pas d'échantillon à collecter par les moyens habituels, et la détection doit reposer entièrement sur la télémétrie d'exécution.

%---------------------------------------------------------------
\subsection{Dissimulation par rootkit : trois couches}

Les rootkits Linux se répartissent selon la couche où ils s'insèrent, chacune offrant un compromis différent entre furtivité et fragilité \cite{elasticRootkits}.

La couche utilisateur repose sur le détournement de l'éditeur de liens dynamique. En déclarant une bibliothèque partagée dans \texttt{/etc/ld.so.preload}, l'attaquant impose que ses propres implémentations soient résolues avant celles de la bibliothèque C standard. Les fonctions de parcours de répertoires ou d'énumération de processus renvoient alors des résultats filtrés, et les outils d'administration usuels deviennent aveugles. La technique est ancienne et triviale à mettre en \oe uvre ; elle reste efficace contre un défenseur qui interroge le système avec les outils du système.

La couche noyau par module chargeable offre une furtivité supérieure au prix d'une contrainte forte : le module doit être compatible avec le noyau cible, et les mécanismes de signature en compliquent le chargement.

La couche eBPF constitue le développement le plus notable des dernières années. Conçue pour permettre l'exécution de programmes vérifiés dans le noyau, cette infrastructure sert désormais aussi à y accrocher des routines de dissimulation. Un rootkit peut intercepter les appels d'énumération de répertoires pour masquer fichiers et processus, et intercepter l'interface de gestion eBPF elle-même pour se soustraire aux outils censés inventorier les programmes chargés.

\paragraph{Un cas instructif.} L'analyse d'une porte dérobée découverte sur une infrastructure cloud compromise, publiée par une équipe de réponse à incident française, documente une architecture en deux modules eBPF, l'un dédié à la dissimulation, l'autre à l'activation du canal de commande sur réception d'un paquet réseau caractéristique \cite{synacktiv2025}. Le point remarquable tient au mécanisme de repli : lorsque la configuration du noyau ne permet pas le fonctionnement des accroches eBPF, le code bascule automatiquement sur la technique de préchargement de bibliothèque décrite plus haut.

Ce comportement mérite d'être souligné, car il documente une capacité d'adaptation à l'environnement d'exécution, obtenue sans aucun recours à l'intelligence artificielle, par une simple logique de repli programmée à l'avance. On y reviendra en section \ref{sec:llm}.

%---------------------------------------------------------------
\subsection{Neutralisation de la télémétrie}

Un agent de sécurité ne voit que ce que le système lui expose. Sous Linux, la télémétrie d'exécution provient principalement du sous-système d'audit du noyau et, pour les outils récents, de sondes eBPF placées sur les points d'entrée des appels système.

Cette dépendance crée une vulnérabilité structurelle. Des travaux récents montrent qu'un code disposant des privilèges suffisants peut agir sur le flux d'événements lui-même, en filtrant ou en supprimant des enregistrements avant qu'ils n'atteignent le tampon lu par l'outil de sécurité. L'agent continue de fonctionner, ses tableaux de bord restent alimentés, et son silence sur une activité donnée est interprété comme une absence d'activité.

La difficulté est ici de nature différente de celle posée par l'obfuscation. Elle ne concerne pas la capacité d'analyse mais la fiabilité de la source, et aucune amélioration du moteur de détection ne la corrige.

%---------------------------------------------------------------
\subsection{Persistance et mimétisme}

Les mécanismes de persistance sous Linux sont nombreux et légitimes, ce qui les rend difficiles à surveiller sans générer de bruit : unités \texttt{systemd}, tâches planifiées, fichiers d'initialisation de shell, préchargement de bibliothèques.

Le mimétisme s'y ajoute. Un service malveillant nommé de manière à évoquer un composant système standard, installé dans une arborescence plausible, ne se distingue d'un service légitime par aucune propriété observable en dehors de son comportement. Le cas cité plus haut recourait précisément à ce procédé. La question n'est donc pas de détecter une anomalie, mais de discriminer entre deux objets dont la seule différence est l'intention.

%---------------------------------------------------------------
\subsection{Commande, contrôle et exfiltration}

Le canal de commande reste le point le plus exposé, et les techniques de dissimulation y visent toutes à réduire la fenêtre d'observation.

L'activation sur paquet caractéristique, dont le cas analysé plus haut fournit une illustration, supprime toute activité réseau tant que le déclencheur n'a pas été reçu : la porte dérobée n'ouvre son port qu'après signal, ce qui la rend invisible à toute cartographie de services. L'usage de services web légitimes comme relais place le trafic dans des flux qu'une organisation ne peut pas bloquer sans coût opérationnel. La génération algorithmique de domaines rend le blocage par liste impraticable.

Une remarque prépare la section suivante : le trafic vers les interfaces de programmation de modèles de langage réunit désormais les propriétés qui font un bon canal de dissimulation, à savoir un volume chiffré, des destinations réputées et un usage professionnel banalisé.

%---------------------------------------------------------------
\subsection{Contrepoint : les spécificités Windows}

Sur Windows, l'évasion à l'exécution s'organise autour de trois surfaces d'observation dont l'équivalent n'existe pas sous Linux. L'interface d'analyse antimalware, qui permet l'inspection des contenus scriptés avant exécution, le mécanisme de traçage d'événements, qui alimente la télémétrie des agents, et l'interception en espace utilisateur, que les agents mettent en place dans les bibliothèques système et que le code malveillant cherche à défaire.

La comparaison éclaire une différence de fond. Sur Windows, l'observation repose largement sur des points d'accroche installés en espace utilisateur, que le code observé peut atteindre. Sous Linux, elle repose sur des primitives noyau, ce qui élève le niveau de privilège requis pour la neutraliser, mais concentre toute la confiance sur une infrastructure unique dont la compromission est d'autant plus lourde de conséquences.

%===============================================================
\section{Intégration de modèles de langage dans la chaîne offensive}
\label{sec:llm}

L'usage de modèles de langage pendant l'exécution d'un code malveillant a été largement commenté depuis 2025. Cette section examine ce que les cas documentés établissent, en séparant les observations des projections.

%---------------------------------------------------------------
\subsection{Les cas documentés}

Un inventaire publié fin 2025 recense les premiers échantillons interrogeant un modèle pendant leur exécution \cite{gtigAI}. Quatre logiques distinctes s'y dégagent, que la vulgarisation confond souvent.

La première est l'auto-réécriture : un chargeur écrit en VBScript interroge périodiquement une interface de modèle pour obtenir de nouvelles variantes obfusquées de son propre code source. C'est le seul cas relevant strictement de cette catégorie, et il vise Windows.

La deuxième est la génération d'instructions à la demande : un outil de collecte attribué à un acteur étatique interroge un modèle hébergé pour produire les commandes à exécuter, sans que son propre code soit modifié.

La troisième détourne les ressources de l'hôte : un voleur de jetons de développement exploite les outils d'IA en ligne de commande déjà installés sur la machine pour rechercher d'autres secrets. Le modèle n'est ici ni embarqué ni distant, mais emprunté.

La quatrième, la plus pertinente pour un environnement Linux, concerne un rançongiciel écrit en Go qui interroge un modèle hébergé localement pour produire à la volée les scripts d'énumération, d'exfiltration et de chiffrement. Présenté d'abord comme le premier rançongiciel assisté par IA \cite{esetPromptLock}, il s'est avéré être le prototype d'un travail académique \cite{nyuRansomware}. Son intérêt tient au modèle local : la détection fondée sur l'observation d'un trafic vers une interface distante y est sans effet, et rien n'interdit ce schéma sur un serveur disposant de ressources de calcul.

%---------------------------------------------------------------
\subsection{Ce que ces cas établissent, et ce qu'ils n'établissent pas}

L'évaluation doit rester mesurée, sous peine de fragiliser le raisonnement défensif qu'elle fonde.

Le chargeur auto-réécrivant est évalué par ses découvreurs comme relevant du développement ou du test, sans capacité de compromettre une cible, et son accès à l'interface du modèle a été révoqué \cite{gtigAI}. Le rançongiciel à modèle local n'a jamais été observé en opération. Plusieurs de ces échantillons étaient détectés par une part significative des moteurs du marché au moment de leur publication.

Ces cas n'établissent donc pas l'existence d'une génération opérationnelle de codes autonomes. Ils établissent trois points plus modestes et plus solides. La faisabilité technique de l'intégration d'un modèle dans la boucle d'exécution est démontrée sur quatre logiques distinctes. Le recours à un modèle local supprime l'indicateur réseau qui constituait la contre-mesure la plus immédiate. Enfin, la trajectoire est ascendante : un rapport ultérieur du même groupe fait état d'un exploit de vulnérabilité inédite dont les analystes estiment avec un haut degré de confiance qu'il a été produit avec l'aide d'un modèle \cite{gtig2026b}.

C'est cette trajectoire, et non l'état des échantillons actuels, qui justifie l'attention défensive.

%---------------------------------------------------------------
\subsection{L'injection de prompt comme technique d'évasion}

Un cas du même inventaire inverse la relation entre modèle et code malveillant. Un shell inversé écrit en PowerShell y intègre des instructions en langage naturel destinées non pas à un modèle utilisé par l'attaquant, mais aux dispositifs d'analyse pilotés par modèle susceptibles d'examiner l'échantillon, auxquels il enjoint de le déclarer inoffensif \cite{gtigAI}.

L'intérêt du cas dépasse sa sophistication, qui est faible. Il montre que confier à un modèle l'évaluation d'un contenu non fiable revient à hériter du problème, non résolu, de l'injection de prompt indirecte. Le point vaut pour tout dispositif de triage qui délègue l'examen d'un fichier suspect à un modèle. Il vaut aussi pour l'analyste : les prompts embarqués dans l'échantillon étudié au chapitre 2 sont du texte écrit par l'attaquant, et rien n'oblige leur auteur à les rédiger à l'intention du seul modèle qu'il interroge.
%---------------------------------------------------------------
\subsection{L'adaptation contextuelle : une capacité déjà présente}

La capacité qui inquiéterait le plus n'est pas la mutation de forme mais l'adaptation à l'environnement : un code qui reconnaîtrait les dispositifs de sécurité en place et ajusterait sa charge en conséquence.

Aucun des échantillons assistés par modèle documentés à ce jour ne présente cette capacité. Mais elle n'est pas hypothétique pour autant, et c'est là que le point de vue Linux change l'analyse : le mécanisme de repli du rootkit eBPF décrit en section \ref{sec:runtime} en est déjà une forme aboutie \cite{synacktiv2025}. Le code y évalue la configuration du noyau et choisit sa technique de dissimulation en conséquence, par une logique de repli programmée à l'avance.

L'adaptation contextuelle est donc une capacité opérationnelle attestée, obtenue sans modèle de langage. Ce que l'intégration d'un modèle changerait n'est pas sa nature mais son étendue : le passage d'un arbre de décision fini, écrit par un développeur et donc énumérable par un analyste, à une génération ouverte dont les branches ne sont pas connues à l'avance. Formulée ainsi, l'hypothèse devient testable, ce qui n'était pas le cas de sa version initiale.

%===============================================================
\section{État de l'art de la détection et de ses limites}
\label{sec:detection}

Quatre approches coexistent, dont aucune ne suffit isolément.

%---------------------------------------------------------------
\subsection{Détection statique par indicateurs}

L'approche ne se réduit pas au condensat cryptographique du fichier. Elle mobilise également des signatures de motifs d'octets, des règles combinant chaînes et conditions structurelles, des empreintes calculées sur des sections ou des tables du binaire, et des fonctions de hachage approché tolérantes aux modifications mineures.

Cette gradation nuance la critique habituelle. Un condensat exact est effectivement mis en échec par la moindre mutation, mais une règle portant sur une routine de déballage ou un hachage approché résistent à une partie du polymorphisme. Ce qui met réellement l'approche en difficulté est l'effet de volume : le coût marginal de production d'une variante inédite est aujourd'hui inférieur au coût marginal de production d'une règle.

Sur ELF, une difficulté supplémentaire tient à la compilation croisée. Une règle validée sur une compilation x86-64 ne se transpose pas mécaniquement aux compilations ARM ou MIPS du même code source, ce qui multiplie l'effort de couverture.

%---------------------------------------------------------------
\subsection{Analyse heuristique et émulation}

L'analyse heuristique recherche des structures de code ou des séquences d'appels caractéristiques sans se fier à une empreinte exacte. Elle s'appuie souvent sur une émulation partielle destinée à franchir les couches de déballage.

Les techniques anti-analyse décrites plus haut la visent directement. Sur Linux, elle se heurte de surcroît à la diversité des architectures, l'émulation devant être disponible et fidèle pour chacune d'elles.

%---------------------------------------------------------------
\subsection{Analyse comportementale}

Les agents de détection observent les actions du programme à l'exécution. Sous Linux, l'outillage contemporain repose massivement sur des sondes eBPF, qui offrent une visibilité fine sur les appels système, les accès fichiers et les connexions réseau, à un coût de performance acceptable.

C'est l'approche la plus robuste face à la variation de forme, puisqu'elle porte sur des invariants d'objectif plutôt que d'implémentation. Ses deux limites ont été exposées plus haut : la source de télémétrie est elle-même attaquable, et un code qui module continûment son comportement en imitant des opérations d'administration légitimes se rapproche progressivement de la distribution normale, ce qui reporte le problème sur la séparation entre dérive légitime et dérive adverse.

%---------------------------------------------------------------
\subsection{Apprentissage automatique}

L'apprentissage automatique s'est imposé pour la classification de binaires, et une revue systématique récente en dresse le panorama \cite{gaber2024}. Trois difficultés conditionnent la valeur de tout dispositif proposé.

La dérive conceptuelle d'abord. Les performances mesurées en validation se dégradent systématiquement en test sur des échantillons postérieurs, effet du vieillissement des modèles face à l'évolution des familles \cite{bodmas}.

Les contraintes opérationnelles ensuite. Un taux de détection annoncé sans le taux de faux positifs associé n'a pas de signification pratique : à l'échelle d'un parc, un taux de faux positifs de 0,1~\% produit un volume d'alertes ingérable. Les protocoles d'évaluation récents imposent pour cette raison une mesure à très faible taux de faux positifs et une séparation temporelle des jeux d'entraînement et de test \cite{ember2024}.

La disponibilité des données enfin, et c'est ici que le périmètre Linux pèse le plus lourd. Les jeux de données publics de référence portent sur des binaires PE ; les corpus ELF disponibles sont plus petits, moins bien étiquetés et souvent restreints aux familles visant les équipements embarqués \cite{crossArchSurvey}. Reproduire sur ELF les protocoles admis dans le monde PE suppose donc de constituer d'abord un corpus. Ce travail n'a pas été mené ici. Le chapitre suivant part d'un seul échantillon, le caractérise, et en tire des règles explicites. Un classifieur entraîné sur ce que j'aurais pu réunir n'aurait pas résisté à l'examen.

%---------------------------------------------------------------
\subsection{Attaques adversariales}

Un classifieur est lui-même une surface d'attaque. La génération d'exemples adverses dans l'espace des problèmes obéit à une contrainte absente du domaine de la vision par ordinateur : la perturbation doit produire un fichier valide, exécutable, et dont la sémantique malveillante est préservée \cite{pierazzi2020}. S'y ajoutent l'empoisonnement des données d'entraînement et l'extraction de modèle.

La condition de recevabilité qu'ils fixent vaut au-delà des classifieurs. Une règle écrite à la main tranche elle aussi entre malveillant et bénin, et une règle validée sur le seul échantillon qui l'a produite ne dit rien de ce qu'elle deviendra face à une variante conçue pour la contourner. Les règles du chapitre 2 n'y échappent pas.

%===============================================================
\section{Synthèse et verrou identifié}

Les défenses actuelles reposent sur une hypothèse de modélisabilité : que l'adversaire produise des artefacts identifiables, des comportements caractérisables, ou à défaut une distribution assez stable pour qu'un classifieur entraîné hier reste valide demain. L'évasion statique attaque la première hypothèse, les techniques d'exécution la deuxième, et la dérive conceptuelle érode la troisième indépendamment de toute intention adverse.

Le contexte Linux accentue chacun de ces effets. La compilation croisée multiplie les formes d'un même comportement. L'exécution sans fichier supprime l'objet même de l'analyse statique. La télémétrie repose sur une infrastructure noyau unique, dont la compromission prive l'agent de toute perception. Et la rareté des corpus publics prive le chercheur des protocoles d'évaluation dont dispose son homologue travaillant sur PE.

L'intégration de modèles de langage dans la chaîne offensive n'introduit pas une menace de nature nouvelle. Les échantillons documentés restent expérimentaux, et l'adaptation à l'environnement, souvent présentée comme leur apport spécifique, est déjà attestée sans eux. Ce que ces modèles modifient est le coût de production de la variation, et l'étendue de l'espace des comportements possibles. Une défense dont le coût marginal d'adaptation reste supérieur au coût marginal de mutation de l'adversaire perd du terrain mécaniquement, quelle que soit sa qualité intrinsèque.